\section{Prospective Tidal Memory Intervention}
\label{app:intervention}
\textbf{Status.} This appendix describes planned work, not an additional
completed experiment. The project contains a local protocol and a one-round
full-memory smoke run; it does not contain completed paired 32-round results
supporting a causal memory claim. No external preregistration identifier is
available in the release. The changes to analysis and measurement specified
below must be frozen before confirmatory runs and are not represented as an
already registered plan.

\subsection{Intervention and Hypothesis}
The 25-agent Tidal Embassy world is initialized once and reused. At round
zero, three agents in the lower archive privately learn that a flood threatens
the only surviving peace treaty unless it is moved and authenticated. The
full-memory condition retains prior social episodes. The ablation removes
prior-round social episodes from decision context while preserving current
locations, inventories, public facts, task state, and directed relationships.
Thus it tests access to episodic social history, not removal of all memory:
tasks, current knowledge, and relationships can also carry historical signals.

The planned seeds are 26080311, 26080312, and 26080313, with 32 rounds per
condition. Initial profiles, world configuration, model and decoding settings,
and activation schedule are matched within each pair. Activation draws should
be precomputed with a separate random stream; reusing a seed alone does not
guarantee matching if the conditions consume randomness differently. API
nondeterminism remains a source of variation.

The hypothesis is that access to earlier social episodes increases warning
diffusion and coordinated rescue within this generated setting. Generative
Agents motivates studying memory in community behavior \cite{park2023};
the separation of information reach and collective action is also informed
by contagion research \cite{centola}.

\begin{figure*}[t]
\centering
\includegraphics[width=0.96\textwidth]{figures/memory_intervention_design.pdf}
\caption{\paneltitle{Prospective design; no paired results reported.}
Matched initial states receive the same private warning. Conditions differ
in access to prior social episodes. Diffusion and material rescue are distinct
outcomes. The planned unit of comparison is a seed pair.}
\label{fig:tidal-design}
\end{figure*}

\subsection{Outcomes and Measurement Controls}
The primary endpoint is warning reach by round 16 among the 22 initially
uninformed agents. Let $Z_i(16)$ indicate that the agent has expressed or
recorded a prespecified treaty-threat proposition by that round. Then
\begin{equation}
D_{16}=\frac{1}{22}\sum_{i\notin I_0}Z_i(16),
\qquad |I_0|=3.
\end{equation}
Measurement must use an append-only audit stream outside the ablated decision
memory; otherwise deleting memory could mechanically reduce measured reach
even after successful transmission. The audit record is not fed back to
agents. Initial warning availability must be preserved equally through its
first intended use, so the ablation does not inadvertently remove the treatment
trigger itself.

Secondary outcomes are first cross-room transmission time, rescue-task
participation and role diversity, treaty relocation and authentication,
unique and repeated dyads, largest component, object transfers, relationship
changes, and rule violations. Unreached transmission or rescue events are
reported as censored at the horizon. Two condition-blind annotators should
label candidate episodes under a frozen proposition codebook; agreement and
resolved disagreements must be reported before condition-level interpretation.

Manipulation checks verify that only the designated agents initially receive
the warning, that paired initial states and activation schedules match, and
that the ablation removes the intended history without exposing hidden facts
through summaries or coordinator feedback. These checks implement the
information-boundary concern highlighted by Zhou et al.\ \cite{misleading}.

\subsection{Analysis and Precision}
The independent comparison unit is a seed pair, not an event or an agent.
For seed $s$, report
$\Delta_s=D_{16,s}^{\mathrm{full}}-D_{16,s}^{\mathrm{ablation}}$, both condition
trajectories, and the material rescue outcome. Report all three differences
and their mean and range. Three pairs are too few for a stable bootstrap
confidence interval or a well-powered confirmatory test; additional seeds
and a prospectively justified sample size are needed for precise inference.

Qualitative selection rules are the earliest cross-room transmission, first
multi-role rescue, largest trust reversal, and first failed coordination.
Selections should be applied identically to both conditions and report when
an event does not occur. The existing one-round smoke run is execution
validation only and must remain outside these comparisons.
